\makeatletter
\let\RGPrefaceChapterHead\@makeschapterhead
\renewcommand{\@makeschapterhead}[1]{\phantomsection\RGPrefaceChapterHead{#1}}
\chapter*{序言}
\let\@makeschapterhead\RGPrefaceChapterHead
\makeatother
\addcontentsline{toc}{chapter}{序言}
\markboth{序言}{序言}
\enlargethispage{\baselineskip}

本书以一门自 2002 年起在 Instituto Superior Técnico (里斯本) 面向数学, 物理和工程专业学生讲授的一学期课程为基础. 本书旨在简要介绍微分几何, 包括微分形式, 随后介绍 Riemann 几何的主要思想 (测地线的极小性质, 完备性和曲率). 最后两章给出若干可能的应用; 这两章也曾分别用于一学期的几何力学和广义相对论课程. 我们希望它们能让数学专业的学生体会 Riemann 几何的用处, 也能让物理和工程专业的学生更有动力学习相关数学背景.

我们假定读者具备线性代数, 多元微积分和微分方程的基本知识, 以及拓扑与代数的初步概念. 为方便读者 (尤其是物理和工程专业的学生), 我们在各章末尾按需归纳了这些背景材料中的主要定义和结果.

为了帮助读者检验并巩固自己的理解, 也为了介绍正文未涉及的重要思想和例子, 我们编入了 330 多道习题, 其中约 140 道在第~\RGref{chap:solutions}~章中给出解答 (全套习题的解答可提供给任课教师). 我们希望本书由此适合自学, 同时仍保留足够多没有给出解答的习题, 供读者挑战.

下面简要介绍各章的内容.

第~\RGref{chap:differentiable-manifolds}~章讨论微分几何的基本概念: 可微流形与可微映射, 向量场和 Lie 括号. 此外, 我们简要介绍 Lie 群和 Lie 群作用.

第~\RGref{chap:differential-forms}~章专门讨论微分形式, 涵盖楔积, 拉回, 外微分, 积分和 Stokes 定理等标准内容.

第~\RGref{chap:riemannian-manifolds}~章引入 Riemann 流形, 讨论 Levi–Civita 联络, 测地线的极小性质和 Hopf–Rinow 定理.

第~\RGref{chap:curvature}~章讨论曲率的概念. 特别地, 我们利用 Cartan 结构方程所提供的有力计算方法证明 Gauss–Bonnet 定理. 我们也讨论常曲率和等距嵌入.

第~\RGref{chap:geometric-mechanics}~章概述几何力学, 包括完整和非完整系统, Lagrange 力学与 Hamilton 力学, 完全可积系统和约化.

第~\RGref{chap:relativity}~章讨论广义相对论, 首先从几何角度介绍狭义相对论. 我们通过 Newton 引力的 Cartan 联络表述引出 Einstein 方程, 并研究 Schwarzschild 解 (包括黑洞) 和宇宙学等基本例子. 最后, 我们讨论因果性以及著名的 Hawking 和 Penrose 奇性定理. 这些定理在入门教材中并不常见, 却是 Riemann 几何很有趣的应用.

最后, 我们感谢许多曾阅读本书全部或部分内容的同事和学生, 感谢他们提出的宝贵意见与建议. 特别感谢我们的同事和朋友 Pedro Girão.
